\documentclass[UTF8,a4paper]{ctexart}
\usepackage{amsmath,amssymb}
\usepackage{graphicx}
\usepackage{geometry}
\usepackage{enumitem}

\usepackage{fancyhdr}
\usepackage{lastpage}


\usepackage{setspace}


\DeclareMathOperator{\arccot}{arccot}
\newcommand{\dlim}{\displaystyle\lim}
\newcommand{\dint}{\displaystyle\int}
\newcommand{\e}{\mathrm{e}}
\renewcommand{\i}{\mathrm{i}}



\geometry{left=31.7mm, right=31.7mm, top=30.4mm, bottom=30.4mm}

\graphicspath{{D:/Git/AmazingEducation/examples/gaokao_sample/resource/gaokao2024/image/}{D:/Git/AmazingEducation/examples/gaokao_sample/resource/gaokao2024/}}

\pagestyle{fancy}
\fancyhead{}
\fancyfoot{}
\fancyfoot[C]{\normalsize 数学试题第\thepage 页（共 \pageref{LastPage} 页）}
\renewcommand{\headrulewidth}{0pt}

\begin{document}


\setstretch{1.5}


\fontsize{11}{13.2}\selectfont



\begin{center}
{\fontsize{17pt}{20.4pt}\selectfont 2024年普通高等学校招生全国统一考试\\[9pt]
\textbf{\fontsize{22pt}{26.4pt}\selectfont 数学}}
\end{center}

\vspace{1em}

\noindent\textbf{注意事项：}\par\smallskip
1. 答卷前，考生务必将自己的姓名、准考证号填写在答题卡上。\par
2. 回答选择题时，选出每小题答案后，用铅笔把答题卡上对应题目的答案标号涂黑。如需改动，用橡皮擦干净后，再选涂其他答案标号。回答非选择题时，将答案写在答题卡上。写在本试卷上无效。\par
3. 考试结束后，将本试卷和答题卡一并交回。

\par\vspace{6pt}



\noindent\textbf{一、选择题：本题共8小题，每小题5分，共40分。在每小题给出的四个选项中，只有一项是符合题目要求的。}\hangindent=2em
\par\medskip




\begin{enumerate}[series=s0, leftmargin=2em, itemsep=0.8em]


\item
已知集合$A=\left\{x\mid-5<x^3<5\right\}$， $B=\left\{-3,-1,0,2,3\right\}$，则$A\cap B=$


\noindent\textbf{A.}\,$\{-1,0\}$\quad\textbf{B.}\,$\{2,3\}$\quad\textbf{C.}\,$\{-3,-1,0\}$\quad\textbf{D.}\,$\{-1,0,2\}$
\par\smallskip





\item
若$\dfrac{z}{z-1}=1+\i$，则$z=$


\noindent\textbf{A.}\,$-1-\i$\quad\textbf{B.}\,$-1+\i$\quad\textbf{C.}\,$1-\i$\quad\textbf{D.}\,$1+\i$
\par\smallskip





\item
已知向量$\boldsymbol a=(0,1)$，$\boldsymbol b=(2,x)$，若 $\boldsymbol b\perp (\boldsymbol b-4\boldsymbol a)$，则 $x=$


\noindent\textbf{A.}\,$-2$\quad\textbf{B.}\,$-1$\quad\textbf{C.}\,$1$\quad\textbf{D.}\,$2$
\par\smallskip





\item
已知$\cos(\alpha+\beta)=m$，$\tan\alpha\tan\beta=2$，则 $\cos(\alpha-\beta)=$


\noindent\textbf{A.}\,$-3m$\quad\textbf{B.}\,$-\dfrac{m}{3}$\quad\textbf{C.}\,$\dfrac{m}{3}$\quad\textbf{D.}\,$3m$
\par\smallskip





\item
已知圆柱和圆锥的底面半径相等，侧面积相等，且它们的高均为$\sqrt 3$，则圆锥的体积为


\noindent\textbf{A.}\,$2\sqrt 3\pi$\quad\textbf{B.}\,$3\sqrt 3\pi$\quad\textbf{C.}\,$6\sqrt 3\pi$\quad\textbf{D.}\,$9\sqrt 3\pi$
\par\smallskip





\item
已知函数$f(x)=\begin{cases}-x^2-2ax-a, &x<0, \\ \e^x+\ln(x+1), &x\geqslant 0\end{cases}$在$\mathbb R$上单调递增，则 $a$ 的取值范围是


\noindent\textbf{A.}\,$(-\infty,0]$\quad\textbf{B.}\,$[-1,0]$\quad\textbf{C.}\,$[-1,1]$\quad\textbf{D.}\,$[0,+\infty)$
\par\smallskip





\item
当$x\in[0,2\pi]$时，曲线 $y=\sin x$与 $y=2\sin\left(3x-\dfrac{\pi}{6}\right)$ 的交点个数为


\noindent\textbf{A.}\,$3$\quad\textbf{B.}\,$4$\quad\textbf{C.}\,$6$\quad\textbf{D.}\,$8$
\par\smallskip





\item
已知函数$f(x)$的定义域为 $\mathbb R$， $f(x)>f(x-1)+f(x-2)$，且当 $x<3$时，$f(x)=x$，则下列结论中一定正确的是


\noindent\textbf{A.}\,$f(10)>100$\quad\textbf{B.}\,$f(20)>1\,000$\quad\textbf{C.}\,$f(10)<1\,000$\quad\textbf{D.}\,$f(20)<10\,000$
\par\smallskip




\end{enumerate}
\vspace{1em}



\noindent\textbf{二、选择题：本题共3小题，每小题6分，共18分。在每小题给出的选项中，有多项符合题目要求。全部选对的得6分，部分选对的得部分分，有选错的得0分。}\hangindent=2em
\par\medskip




\begin{enumerate}[series=s1, leftmargin=2em, itemsep=0.8em]


\item
为了解推动出口后的亩收入（单位：万元）情况，从该种植区抽取样本，得到推动出口后亩收入的样本均值$\overline{x}=2.1$， 样本方差$s^2=0.01$，已知该种植区以往的亩收入$X$服从正态分布$N\left(1.8, 0.1^2\right)$，假设推动出口后的亩收入$Y$服从正态分布$N\left(\overline{x},s^2\right)$，则（若随机变量$Z$服从正态分布$N\left(\mu,\sigma^2\right)$，则$P(Z<\mu+\sigma)\approx 0.841\,3$）


\noindent\textbf{A.}\,$P(X>2)>0.2$\quad\textbf{B.}\,$P(X>2)<0.5$\quad\textbf{C.}\,$P(Y>2)>0.5$\quad\textbf{D.}\,$P(Y>2)<0.8$
\par\smallskip





\item
设函数$f(x)=\left(x-1\right)^2\left(x-4\right)$，则




\begin{tabular}{@{}p{0.47\linewidth}@{\hspace{0.06\linewidth}}p{0.47\linewidth}@{}}
\textbf{A.}\,$x=3$是$f(x)$的极小值点 & \textbf{B.}\,当$0<x<1$时，$f(x)<f\left(x^2\right)$ \\
\textbf{C.}\,当$1<x<2$时，$-4<f(2x-1)<0$ & \textbf{D.}\,当$-1<x<0$时，$f(2-x)>f(x)$ \\
\end{tabular}
\par\smallskip






\item
已知曲线$C$的一部分可看作某种丝带造型（原卷配图略）．已知$C$过坐标原点$O$，且$C$上的点满足横坐标大于$-2$，到点$F(2,0)$的距离与到定直线$x=a$ $(a < 0)$的距离之积为$4$，则




\begin{tabular}{@{}p{0.47\linewidth}@{\hspace{0.06\linewidth}}p{0.47\linewidth}@{}}
\textbf{A.}\,$a=-2$ & \textbf{B.}\,点$\left(2\sqrt 2,0\right)$在$C$上 \\
\textbf{C.}\,$C$在第一象限的点的纵坐标的最大值为$1$ & \textbf{D.}\,当点$(x_0,y_0)$在$C$上时， $y_0\leqslant \dfrac{4}{x_0+2}$ \\
\end{tabular}
\par\smallskip





\end{enumerate}
\vspace{1em}



\noindent\textbf{三、填空题：本题共3小题，每小题5分，共15分。}\hangindent=2em
\par\medskip




\begin{enumerate}[series=s2, leftmargin=2em, itemsep=0.8em]


\item
设双曲线$C:\dfrac{x^2}{a^2}-\dfrac{y^2}{b^2}=1$ $(a>0,b>0)$的左、右焦点分别为$F_1$，$F_2$，过$F_2$作平行于$y$轴的直线交 $C$于 $A$、$B$两点，若$\left|F_1A\right|=13$，$\left|AB\right|=10$，则$C$的离心率为\underline{\qquad\qquad}．




\item
若曲线$y=\e^x+x$在点$(0,1)$处的切线也是曲线$y=\ln(x+1)+a$的切线，则$a=$\underline{\qquad\qquad}．




\item
甲乙两人各有四张卡片，每张卡片上标有一个数字，甲的卡片上分别标有数字$1$，$3$，$5$，$7$，乙的卡片上分别标有数字$2$，$4$，$6$，$8$．两人进行四轮比赛，在每轮比赛中，两人各自从自己持有的卡片中随机选一张，并比较所选卡片上的数字大小，数字大的人得$1$分，数字小的人得$0$分，然后各自弃置此轮所选的卡片（弃置的卡片在此后的轮次中不能使用），则四轮比赛后，甲的总得分不小于$2$的概率为\underline{\qquad\qquad}．



\end{enumerate}
\vspace{1em}



\noindent\textbf{四、解答题：本题共5小题，共77分。解答应写出文字说明、证明过程或演算步骤。}\hangindent=2em
\par\medskip




\begin{enumerate}[series=s3, leftmargin=2em, itemsep=0.8em]


\item
设$\triangle ABC$的内角 $A$，$B$，$C$的对边分别为 $a$， $b$， $c$，已知 $\sin C=\sqrt 2\cos B$， $a^2+b^2-c^2=\sqrt 2 ab$．

（1）求$B$；

（2）若$\triangle ABC$的面积为 $3+\sqrt 3$，求 $c$．





\item
已知$A(0,3)$和 $P\left(3,\dfrac{3}{2}\right)$为椭圆$\dfrac{x^2}{a^2}+\dfrac{y^2}{b^2}=1$ $(a>b>0)$上两点．

（1）求$C$的离心率；

（2）若过$P$的直线 $l$交 $C$于另一点$B$，且 $\triangle ABP$的面积为 $9$，求 $l$的方程．





\item
四棱锥$P$--$ABCD$中， $PA\perp$底面$ABCD$， $PA=AC=2$， $BC=1$， $AB=\sqrt{3}$（原卷立体几何配图略）．

（1）若$AD\perp PB$，证明： $AD\parallel$平面$PBC$；

（2）若$AD\perp DC$，且二面角 $A$--$CP$--$D$的正弦值为 $\dfrac{\sqrt{42}}{7}$，求$AD$．





\item
已知函数$f(x)=\ln \dfrac{x}{2-x}+ax+b\left(x-1\right)^3$．

（1）若$b=0$，且 $f'(x)\geqslant 0$，求 $a$的最小值；

（2）证明：曲线$y=f(x)$是中心对称图形；

（3）若 $f(x)>-2$当且仅当 $1<x<2$，求 $b$的取值范围．





\item
设$m$为正整数，数列$a_1$，$a_2$，$\cdots$，$a_{4m+2}$是公差不为$0$的等差数列，若从中删去两项$a_i$和$a_j$ $(i<j)$ 后剩余的$4m$项可被平均分为$m$组，且每组的$4$个数都能构成等差数列，则称数列$a_1$，$a_2$，$\cdots$，$a_{4m+2}$是$(i,j)$--可分数列．

（1）写出所有的$(i,j)$，$1\leqslant i<j\leqslant 6$，使得数列$a_1$，$a_2$，$\cdots$，$a_{6}$是$(i,j)$--可分数列；

（2）当$m\geqslant 3$时，证明：数列$a_1$，$a_2$，$\cdots$，$a_{4m+2}$是$(2,13)$--可分数列；

（3）从$1$，$2$，$\cdots$，$4m+2$中一次任取两个数$i$和$j$ $(i<j)$，记数列$a_1$，$a_2$，$\cdots$，$a_{4m+2}$是$(i,j)$--可分数列的概率为$P_m$，证明：$P_m>\dfrac{1}{8}$．




\end{enumerate}
\vspace{1em}


\end{document}